% ============================================================
%  Trigonometria: Adição de Arcos — Apostila Premium | PratiqueJá
%  Engine: XeLaTeX
% ============================================================
\documentclass[12pt]{article}
\usepackage{../../pratiqueja}

\newcommand{\hlm}[1]{{\color{darkblue}#1}}

\setsubject{Trigonometria: Adição de Arcos}

\begin{document}
\pagenumbering{arabic}
\setstretch{1.2}
\color{bodytext}

\heroheader{Trigonometria: Adição de Arcos}%
           {Fórmulas de adição, duplicação, bissecção e produto}%
           {Matemática · Avançado}

\setlength{\columnsep}{18pt}
\setlength{\columnseprule}{0.5pt}
\def\columnseprulecolor{\color{exborder}}

\begin{multicols}{2}

\TW{Def}{darkblue}{Fórmulas de Adição e Subtração}{Seno, cosseno e tangente da soma/diferença de arcos}

\regra{
\textbf{Seno:}
\[
  \sin(a + b) = \sin a \cos b + \cos a \sin b
\]
\[
  \sin(a - b) = \sin a \cos b - \cos a \sin b
\]
\textbf{Cosseno:}
\[
  \cos(a + b) = \cos a \cos b - \sin a \sin b
\]
\[
  \cos(a - b) = \cos a \cos b + \sin a \sin b
\]
\textbf{Tangente:}
\[
  \tan(a + b) = \frac{\tan a + \tan b}{1 - \tan a \tan b}
\]
\[
  \tan(a - b) = \frac{\tan a - \tan b}{1 + \tan a \tan b}
\]
(válidas quando os denominadores $\neq 0$)
}

\exemp{
\textbf{Exemplo 1 — $\sin 75°$:}

$\sin 75° = \sin(45° + 30°)$

$= \sin 45°\cos 30° + \cos 45°\sin 30°$

$= \dfrac{\sqrt{2}}{2}\cdot\dfrac{\sqrt{3}}{2} + \dfrac{\sqrt{2}}{2}\cdot\dfrac{1}{2}$

$\hlm{= \dfrac{\sqrt{6}+\sqrt{2}}{4}}$

\textbf{Exemplo 2 — $\cos 15°$:}

$\cos 15° = \cos(45° - 30°)$

$= \cos 45°\cos 30° + \sin 45°\sin 30°$

$= \dfrac{\sqrt{2}}{2}\cdot\dfrac{\sqrt{3}}{2} + \dfrac{\sqrt{2}}{2}\cdot\dfrac{1}{2}$

$\hlm{= \dfrac{\sqrt{6}+\sqrt{2}}{4}}$
}

\TW{2a}{darkblue}{Fórmulas de Duplicação (Arco Duplo)}{Obtidas fazendo $b = a$ nas fórmulas de adição}

\regra{
\[
  \sin(2a) = 2\sin a \cos a
\]
\[
  \cos(2a) = \cos^2 a - \sin^2 a
\]
\[
  = 1 - 2\sin^2 a = 2\cos^2 a - 1
\]
\[
  \tan(2a) = \frac{2\tan a}{1 - \tan^2 a}
\]
}

\exemp{
\textbf{Exemplo 3 — $\sin(2a)$ com $\sin a = \frac{3}{5}$, 1º quadrante:}

$\cos a = \dfrac{4}{5}$ (Pitágoras)

$\sin(2a) = 2 \cdot \dfrac{3}{5} \cdot \dfrac{4}{5} = \hlm{\dfrac{24}{25}}$

\textbf{Exemplo 4 — $\cos^2 a - \sin^2 a$ para $a = 30°$:}

$\cos(60°) = \dfrac{1}{2}$

Verif.: $\left(\dfrac{\sqrt{3}}{2}\right)^2 - \left(\dfrac{1}{2}\right)^2 = \dfrac{3}{4} - \dfrac{1}{4} = \hlm{\dfrac{1}{2}}$ \;\ok
}

\TW{a/2}{badgepurple}{Fórmulas de Bissecção (Arco Metade)}{Expressam $\sin^2$ e $\cos^2$ em termos de $\cos(2a)$}

\regra{
Das fórmulas de duplicação do cosseno:
\[
  \sin^2 a = \frac{1 - \cos(2a)}{2} \qquad \cos^2 a = \frac{1 + \cos(2a)}{2}
\]
Fazendo $a \to \frac{a}{2}$:
\[
  \sin^2\!\tfrac{a}{2} = \frac{1 - \cos a}{2} \qquad \cos^2\!\tfrac{a}{2} = \frac{1 + \cos a}{2}
\]
}

\exemp{
\textbf{Exemplo 5 — $\cos^2 15°$:}

$\cos^2 15° = \dfrac{1 + \cos 30°}{2} = \dfrac{1 + \frac{\sqrt{3}}{2}}{2}$

$\hlm{= \dfrac{2 + \sqrt{3}}{4}}$
}

\TW{P→S}{badgepurple}{Produto em Soma (e Soma em Produto)}{Simplificam expressões trigonométricas}

\regra{
\textbf{Produto em soma:}
\[
  \sin a \cos b = \tfrac{1}{2}[\sin(a+b) + \sin(a-b)]
\]
\[
  \cos a \cos b = \tfrac{1}{2}[\cos(a-b) + \cos(a+b)]
\]
\[
  \sin a \sin b = \tfrac{1}{2}[\cos(a-b) - \cos(a+b)]
\]

\textbf{Soma em produto} (com $p = a+b$, $q = a-b$):
\[
  \sin p + \sin q = 2\sin\!\tfrac{p+q}{2}\cos\!\tfrac{p-q}{2}
\]
\[
  \cos p + \cos q = 2\cos\!\tfrac{p+q}{2}\cos\!\tfrac{p-q}{2}
\]
}

\exemp{
\textbf{Exemplo 6 — simplificar $\sin 70° + \sin 10°$:}

$= 2\sin\!\dfrac{80°}{2}\cos\!\dfrac{60°}{2} = 2\sin 40°\cos 30°$

$\hlm{= \sqrt{3}\sin 40°}$

\textbf{Exemplo 7 — provar $\cos^2 a - \sin^2 b = \cos(a+b)\cos(a-b)$:}

Lado direito:

$= [\cos a\cos b - \sin a\sin b][\cos a\cos b + \sin a\sin b]$

$= \cos^2 a\cos^2 b - \sin^2 a\sin^2 b$

$= \cos^2 a(1{-}\sin^2 b) - \sin^2 b(1{-}\cos^2 a)$

$\hlm{= \cos^2 a - \sin^2 b}$ \;\ok
}

\nota{As fórmulas de adição são a base de todas as demais: duplicação, bissecção e produto em soma derivam diretamente delas.}

\TW{Ref}{badgepurple}{Valores de Referência}{Ângulos notáveis de 0° a 90°}

\regra{
\begin{center}
\footnotesize
\renewcommand{\arraystretch}{1.7}
\begin{tabular}{|c|c|c|c|}
\hline
$\theta$ & $\sin\theta$ & $\cos\theta$ & $\tan\theta$ \\
\hline
$0°$ & $0$ & $1$ & $0$ \\
\hline
$30°$ & $\frac{1}{2}$ & $\frac{\sqrt{3}}{2}$ & $\frac{\sqrt{3}}{3}$ \\
\hline
$45°$ & $\frac{\sqrt{2}}{2}$ & $\frac{\sqrt{2}}{2}$ & $1$ \\
\hline
$60°$ & $\frac{\sqrt{3}}{2}$ & $\frac{1}{2}$ & $\sqrt{3}$ \\
\hline
$90°$ & $1$ & $0$ & indef. \\
\hline
\end{tabular}
\end{center}
}

\end{multicols}

\end{document}
